\section{Представление о языках программирования высокого уровня}

\textbf{Язык программирования} --- формальный язык для записи алгоритмов и структур данных в форме, допускающей автоматическое выполнение после трансляции или интерпретации. \textbf{Языки высокого уровня} предоставляют абстракции, в большей степени ориентированные на алгоритм и предметную область, чем на систему команд конкретного процессора.

Примеры: Fortran, C, C++, Java, Python, C\#, Julia. Граница между ``высоким'' и ``низким'' уровнем условна: например, C и C++ являются языками высокого уровня, но дают программисту сравнительно прямой контроль над памятью.

\subsection{Синтаксис и семантика}

\textbf{Синтаксис} определяет правила построения корректных конструкций языка; \textbf{семантика} --- их смысл и результат выполнения. Например, в императивном языке
\[
x=x+1
\]
означает не математическое равенство, а присваивание нового значения переменной.

Формальный язык включает лексические элементы, выражения, объявления, операторы управления и правила организации модулей. Синтаксис часто описывают формальной грамматикой, например формой Бэкуса--Наура.

\subsection{Типы данных и управление вычислением}

\textbf{Тип} задаёт множество допустимых значений и операций над ними. К базовым и составным типам относятся целые и вещественные числа, логические значения, строки, массивы, записи/структуры, указатели и ссылки, классы и пользовательские типы.

При \textbf{статической типизации} основные типы выражений известны и проверяются до выполнения программы; при \textbf{динамической типизации} тип значения определяется во время выполнения.

Для структурного программирования центральны три управляющие конструкции:
\begin{itemize}
\item следование;
\item ветвление;
\item цикл.
\end{itemize}
Из них вместе с функциями, процедурами и модулями строятся более сложные алгоритмы.

\subsection{Основные парадигмы}

Современные языки часто мультипарадигменны.

\begin{itemize}
\item \textbf{Императивное/процедурное программирование}: программа --- последовательность операций, изменяющих состояние; задача декомпозируется на процедуры и функции.
\item \textbf{Объектно-ориентированное}: программа строится из объектов и классов; основные механизмы --- инкапсуляция, наследование и полиморфизм.
\item \textbf{Функциональное}: вычисления выражаются через композицию функций, функции высших порядков и по возможности неизменяемые данные.
\item \textbf{Логическое и декларативное}: задаются отношения, правила или желаемый результат, а способ поиска решения в значительной степени определяется системой выполнения.
\end{itemize}

\subsection{Компиляция и интерпретация}

Для выполнения исходный текст должен быть преобразован в исполняемую форму.

\textbf{Компилятор} анализирует программу и переводит её в машинный код или промежуточное представление. Типичная цепочка:
\[
\text{лексический анализ}
\rightarrow
\text{синтаксический анализ}
\rightarrow
\text{семантический анализ}
\rightarrow
\text{оптимизация}
\rightarrow
\text{генерация кода}.
\]

\textbf{Интерпретатор} организует выполнение программы без предварительного получения отдельного машинного исполняемого файла в традиционном смысле. На практике распространены гибридные схемы: байткод, виртуальные машины и JIT-компиляция.

Разница между ``компилируемым'' и ``интерпретируемым'' языком относится скорее к реализации, чем к математической сущности языка: один и тот же язык может иметь разные системы исполнения.

\subsection{Абстракции высокого уровня и библиотеки}

Высокоуровневые языки освобождают программиста от непосредственной записи машинных команд и предоставляют:
\begin{itemize}
\item функции, модули и пространства имён;
\item абстрактные типы и классы;
\item стандартные контейнеры и алгоритмы;
\item механизмы обработки ошибок;
\item библиотеки ввода-вывода, линейной алгебры, сетей, графики и т.д.
\end{itemize}

В научных вычислениях существенны численная производительность, поддержка массивов, библиотеки линейной алгебры, средства параллельных вычислений и взаимодействие с кодом на других языках. Поэтому выбор языка определяется не ``лучшим языком вообще'', а классом решаемых задач.


\subsection{Формальный язык: лексика, синтаксис и семантика}

Программа является текстом формального языка. На лексическом уровне она
разбивается на токены: идентификаторы, литералы, ключевые слова, знаки операций.
Синтаксис определяет допустимую структуру конструкций, например грамматикой
в форме Бэкуса--Наура. Семантика задаёт смысл синтаксически корректной
программы.

Поэтому ``синтаксически правильно'' не означает ``семантически корректно''.
Например, выражение может соответствовать грамматике, но нарушать правила типов
или обращаться к несуществующему имени.

\subsection{Статическая и динамическая типизация}

\textbf{Система типов} классифицирует значения и операции. При статической
типизации значительная часть проверок выполняется до запуска программы; при
динамической --- тип значения и корректность операции могут проверяться во
время выполнения.

Различают также явное и выводимое задание типов. Параметрический полиморфизм
позволяет описывать алгоритмы над семейством типов, например
\verb|Vector<T>| или обобщённую функцию сортировки.

Типы служат не только представлению данных, но и средством спецификации:
хорошо выбранный тип способен исключить целый класс некорректных состояний ещё
до выполнения программы.

\subsection{Этапы трансляции программы}

Типичная компиляция имеет вид
\[
\begin{aligned}
\text{текст}&\rightarrow\text{лексический анализ}
\rightarrow\text{синтаксический анализ}
\rightarrow\text{семантический анализ}\\
&\rightarrow\text{промежуточное представление}
\rightarrow\text{оптимизация}
\rightarrow\text{машинный код}.
\end{aligned}
\]
После компиляции объектные модули и библиотеки связываются компоновщиком.
Интерпретатор выполняет программу или её промежуточное представление во время
работы; JIT-компиляция сочетает оба подхода, переводя ``горячие'' участки в
машинный код динамически.

Язык сам по себе не обязан быть строго ``компилируемым'' или
``интерпретируемым'': это свойство конкретной реализации.

\subsection{Среда выполнения и управление памятью}

Во время выполнения используются стек вызовов, статическая область и динамическая
память. Автоматическое управление памятью может осуществляться сборщиком мусора;
в системных языках время жизни ресурсов часто контролируется программистом или
системой владения.

Для научных вычислений выбор языка связан не только с удобством синтаксиса, но и
с моделью данных, стоимостью абстракций, доступом к векторизации/параллелизму и
качеством библиотек линейной алгебры, сеток и визуализации.



\subsection{Модульность, библиотеки и взаимодействие языков}

Язык высокого уровня обычно используется не изолированно, а вместе с
библиотеками и системой модулей. Модуль задаёт пространство имён и публичный
интерфейс, скрывая детали реализации. Это уменьшает число глобальных связей и
позволяет независимо компилировать и тестировать части программы.

Научные приложения часто являются многоязыковыми: управляющая логика может быть
на Python, вычислительное ядро --- на C/C++ или Fortran, параллельные ядра --- на
CUDA/OpenCL. Связь обеспечивается ABI и FFI. Поэтому при выборе языка важна не
только его собственная скорость, но и возможность без лишних копирований
работать с существующими библиотеками и форматами данных.

\subsection{Критерии выбора языка}

Для конкретной задачи разумно сравнивать:
\begin{itemize}
    \item выразительность и вероятность программных ошибок;
    \item качество компилятора и производительность;
    \item модель памяти и средства безопасного управления ресурсами;
    \item библиотеки численных методов, ввода-вывода и визуализации;
    \item поддержку многопоточности, SIMD, распределённых и GPU-вычислений;
    \item переносимость между ОС и архитектурами;
    \item средства тестирования, профилирования и отладки;
    \item зрелость экосистемы и стоимость сопровождения кода.
\end{itemize}

Поэтому разделение на ``быстрые'' и ``медленные'' языки слишком грубо:
производительность определяется конкретной реализацией, алгоритмом, библиотеками
и долей времени, проведённой в оптимизированном вычислительном ядре.


\subsection{Что важно сказать на экзамене}

Нужно определить язык высокого уровня и противопоставить его машинному коду/ассемблеру, затем объяснить:
\[
\boxed{\text{синтаксис и семантика}
\rightarrow
\text{типы данных}
\rightarrow
\text{парадигмы}
\rightarrow
\text{компиляция/интерпретация}}.
\]
Желательно привести примеры языков и отметить, что уровень языка означает прежде всего степень абстрагирования от конкретной архитектуры ЭВМ.

\newpage
